\section{Related literature and what is different here}\label{sec:lit}

The loss ratio touches four literatures: the axiomatic theory of premia and risk measures, the
theory of performance ratios and acceptability indices, implicitly defined risk functionals, and
ratio representations in decision theory and statistics. For each we say what is known, what we use,
and where our problem differs. Section~\ref{sec:lit-table} collects the differences in a table.

\subsection{Distortion premia, Choquet integrals and signed Choquet integrals}

The axiomatic theory of insurance premia starts from properties of a single premium functional
$H$ on losses: translation invariance, $H(X+c)=H(X)+c$, and ``no unjustified loading'', $H(c)=c$ for constant
losses \citep{gerber1974},
monotonicity, law invariance, and, most importantly, additivity for comonotonic risks. The dual
theory of \citet{yaari1987} replaces utility transformations of payoffs by distortions of probabilities, which
leads to the Choquet-type representation $\int g(P(X>x))\,dx$ and to additivity for comonotonic risks, and \citet{wang1996} introduced it as a premium principle. \citet{wyp1997}
characterize the \emph{distortion premium} $C_g$ by law invariance, monotonicity, comonotonic
additivity, a continuity condition and $H(1)=1$. The mathematical backbone is Schmeidler's
(\citeyear{schmeidler1986}) theorem: a monotone, comonotonic additive functional on bounded
measurable functions is a Choquet integral with respect to a capacity, with no continuity needed;
see also \citet{greco1982} and \citet{denneberg1994} on cones. Coherent risk measures are defined by
monotonicity, translation invariance, subadditivity and positive homogeneity \citep{adeh1999,delbaen2002}; in the
law-invariant case, comonotonic coherent risk measures are mixtures of expected shortfalls over levels in $[0,1]$
\citep{kusuoka2001,jst2006,folmerschied2016}.

Without monotonicity the corresponding representation theory is that of \emph{signed} Choquet
integrals. \citet{msm1994} and \citet{cmmm2012,cmmm2015} show that comonotonic additive functionals
of bounded variation are exactly differences of Choquet integrals with respect to set functions of
bounded variation. \citet{www2020} prove the law-invariant version used here: on an atomless
$L^\infty$, a functional is law invariant, comonotonic additive and uniformly norm continuous if and
only if it is a signed Choquet integral $I_h$ with $h$ of bounded variation. They also show that
$I_h$ is convex if and only if $h$ is concave, and \citet{wangwei2020} study when level sets of such
functionals are convex. \citet{www2020b} extend the characterization to law-invariant convex cones containing
$L^\infty$, such as $L^p$, by adding continuity at infinity, which holds automatically on $L^\infty$, and give a
representation of convex signed Choquet integrals by mixtures of expected shortfalls.

These theorems underlie our representation step, and
Theorems~\ref{thm:A} and~\ref{thm:B} are the loss-ratio counterparts of \citet[Thm.~1]{www2020} and
of \citet{wyp1997}. We re-prove the representation directly (Steps~5--8 of Section~\ref{sec:proofs})
to keep track of constants and because we need no continuity in the monotone tier.

All of these results characterize one functional by axioms placed on that
functional. A loss ratio is a ratio of two, and the axioms must be placed on the ratio, which does
not determine its parts (Section~\ref{sec:intro}). Bridging this gap is the content of
Sections~\ref{sec:axioms} and~\ref{sec:proofs}. We are not aware of an if-and-only-if
characterization of a ratio of two Choquet integrals, or of two risk measures, in this literature.

\subsection{Performance ratios and acceptability indices}

The best-known ratios of risk functionals are performance measures. The gain--loss ratio of
\citet{bernardoledoit2000} divides the expected positive part of a payoff by the expected negative
part, and \citet{biaginipinar2013} study its optimization. \citet{chernymadan2009} give an
axiomatic theory of \emph{acceptability indices}: a map $\alpha$ that is monotone, quasi-concave,
scale invariant and Fatou is characterized by an increasing family of coherent risk measures
$\{\rho_x\}$ with $\alpha(X)=\sup\{x:\rho_x(X)\le0\}$, and the family is recovered by the cash-shift
formula $\rho_x(X)=\inf\{c:\alpha(X+c)\ge x\}$. Dynamic versions and the correspondence between
indices and families are developed by \citet{bcz2014} and \citet{bbn2014}, and the correspondence
between quasi-convex functionals and acceptance families by \citet{drapeaukupper2013}. The
Sharpe ratio is not an index in this sense, and the conic-finance programme of
\citet{chernymadan2010} uses the \emph{spread} between ask and bid prices instead of a ratio.
\citet{ck2013} propose reward--risk ratios, a reward functional divided by a risk functional, and
list axioms that make such a ratio a performance measure, giving sufficient conditions rather than a
characterization. 

\citet{hlw2025} give an axiomatic theory of diversification indices. They note that the diversification
ratio $\phi(\sum_iX_i)/\sum_i\phi(X_i)$ of a risk measure $\phi$ may be negative, needs the conventions $0/0=0$
and $1/0=\infty$, and is not location invariant, and they characterize instead the \emph{diversification
quotient} $\alpha^*/\alpha$, where $\alpha^*$ is the level at which a decreasing family $(\rho_\beta)$ of risk
measures, evaluated at the pooled portfolio, reaches $\sum_i\rho_\alpha(X_i)$. Six axioms (non-negativity,
location and scale invariance, rationality, normalization and continuity) characterize the quotients built from
monotone, cash-additive and positively homogeneous families \citep[Thm.~1]{hlw2025}, and portfolio convexity
restricts the class to coherent families when there are at least four assets \citep[Thm.~2]{hlw2025}; they also
relate the quotient to the acceptability indices of \citet{chernymadan2009}. This is the closest if-and-only-if
result of association type. It shares two features with ours: a ratio-type quantity is recovered as the level at
which a family of risk measures crosses a threshold (compare Corollary~\ref{cor:assoc}), and the naive ratio
needs conventions where its denominator vanishes (compare the labels $\infty$ and $\star$). It differs in the
object, an index on vectors of risks that is invariant under location and scale, whereas a loss ratio is a
function of a single book that moves along each cash line as a M\"obius map; and in monotonicity, which their
rationality axiom uses and which a loss ratio cannot have (Proposition~\ref{prop:F2}).

The cash-shift recovery formula of \citet{chernymadan2009} and its
generalizations is the model for our reading of $\rho$ and $\Pi$ from cash lines, and the association is the
loss-ratio analogue of the level-crossing definition of an index. Three differences matter.

First, \emph{direction}. In an acceptability index, larger gains are better, so
the index is monotone and the family $\{\rho_x\}$ is monotone in the level. In a reward--risk ratio,
the numerator is a reward and the denominator a risk, so the ratio is monotone again. In a loss
ratio, numerator and denominator are both increasing in the loss. We show that the ratio is then
not monotone (Proposition~\ref{prop:F2}), and that the associated family
$\rho_\beta=(\rho-\beta\Pi)/(1-\beta)$ is not monotone either (Corollary~\ref{cor:assoc}): it is a
\emph{signed} Choquet integral even when $\rho$ and $\Pi$ are distortion premia. The index theorems
therefore do not apply, and the signed tier of Theorem~\ref{thm:A} is not a luxury but the natural
home of the associated family. Second, \emph{cardinality}. Indices are ordinal objects, determined up to
increasing transformations, whereas the loss ratio is cardinal: it obeys an accounting identity
under pooling (Section~\ref{sec:cm}). Third, \emph{what is characterized}. We characterize
the ratio together with its numerator and denominator, not an index through a family.

\subsection{Level-crossing and implicitly defined functionals}

Several risk functionals are defined implicitly, by where a monotone family crosses a level. Orlicz
premia solve $\E[\varphi(X/\pi)]=1-\alpha$ for $\pi$ \citep{hg1982,brg2008}, and PELVE is the
level at which expected shortfall equals a given value-at-risk
\citep{liwang2023,fiorirg2023}. Return risk measures \citep{blrg2018} and star-shaped risk measures
\citep{castagnoli2022} replace cash additivity by positive homogeneity or star-shapedness. The
association belongs to this family: it defines $\LR(X)$ as the level $\beta$ at which a
cash-additive family $\rho_\beta$ vanishes.

The results in this literature establish \emph{existence, uniqueness and
comparative statics} of the implicit functional. They do not represent it as a ratio of two
explicit functionals. In our setting the implicit definition is recovered as a corollary
(Corollary~\ref{cor:assoc}), and it holds exactly off the set $\{\rho=\Pi=0\}$, not on all of $L^\infty$,
a point that the implicit definition alone obscures.

\subsection{Ratio representations in decision theory and statistics}

In decision theory, \citet{chew1983} shows that weighted-utility preferences are represented by a
\emph{ratio} of two linear functionals, while \citet{dekel1986} and \citet{fishburn1983} show that
the weaker \emph{betweenness} property yields only an implicit, ordinal representation. This
distinction is exactly the one we meet: a mediant inequality for pooled loss ratios is a betweenness
property and cannot characterize the ratio, whereas the premium-weighted identity is the
weighted-utility analogue and can (Section~\ref{sec:cm} and the cube witness in
Section~\ref{sec:necessity}). In statistics, \citet{gneiting2011} shows that ratios of expectations
$\E[r(Y)]/\E[s(Y)]$ are elicitable, which is the only published if-and-only-if statement about a
ratio functional that we know of; it concerns forecasting, not pricing, and has no counterpart of
pooling or cash shifts.

\subsection{Summary of the differences}\label{sec:lit-table}

\begin{table}[ht]
\centering\small
\begin{tabularx}{\textwidth}{@{}>{\raggedright\arraybackslash}p{3.1cm}>{\raggedright\arraybackslash}X>{\raggedright\arraybackslash}X@{}}
\toprule
\textbf{Literature} & \textbf{What it characterizes} & \textbf{Difference from this paper}\\
\midrule
Yaari (1987); Wang, Young, Panjer (1997)
 & One premium functional $H$; axioms on $H$ (law invariance, monotonicity, comonotonic additivity).
 & We characterize a \emph{ratio} of two such functionals with axioms on the ratio. Monotonicity is placed on the \emph{implied premium and risk measure}, never on the ratio.\\
Schmeidler (1986); Murofushi--Sugeno--Machida (1994); Cerreia-Vioglio et al.\ (2012, 2015); Wang--Wei--Willmot (2020)
 & Representation of one comonotonic additive functional as a (signed) Choquet integral.
 & Used as the engine (applied to $\rho$ and $\Pi$). New: the axioms are on $\LR$, comonotonic additivity of $\rho$ and $\Pi$ is \emph{derived} (also on the break-even set), and uniqueness is unconditional.\\
Cherny--Madan (2009); Bielecki--Cialenco--Zhang (2014); Biagini--Bion-Nadal (2014)
 & Monotone, quasi-concave indices $\Leftrightarrow$ increasing families of coherent measures.
 & $\LR$ cannot be monotone (Prop.~\ref{prop:F2}); the associated family is signed and non-monotone. Cardinal (pooling identity) rather than ordinal.\\
Cheridito--Kromer (2013)
 & Axioms for reward--risk ratios (sufficient conditions).
 & Reward over risk is monotone; loss over premium is not. We give an \emph{iff}, within a stated class.\\
Han--Lin--Wang (2025)
 & Diversification quotient: iff axiomatization of a level-crossing index.
 & Closest in spirit; monotone families and a portfolio index. We characterize the ratio and its parts.\\
Haezendonck--Goovaerts (1982); Li--Wang (2023); Bellini et al.\ (2018)
 & Implicitly defined premia/risk measures, return risk measures.
 & Existence and uniqueness of an implicit functional, not a ratio representation. The association is recovered as a corollary and shown to hold only off $\{\rho=\Pi=0\}$.\\
Chew (1983); Dekel (1986)
 & Ratio (weighted utility) versus implicit (betweenness) representations.
 & Explains why a mediant inequality is not enough and a premium-weighted identity is needed.\\
Gneiting (2011)
 & Elicitability of ratios of expectations.
 & Statistical, no axioms on pooling or cash shifts.\\
\bottomrule
\end{tabularx}
\caption{Positioning of the paper. ``Iff'' means a characterization by necessary and sufficient axioms.}
\end{table}

None of our proofs depends on a citation: every load-bearing step is proved in the text.
